%---------------------------------------------------------------------------%
%->> Section 2.3: 传统微服务故障诊断方法
%---------------------------------------------------------------------------%

\section{传统微服务故障诊断方法}

传统微服务故障根因定位方法大致可分为两类：因果推断的方法和特征学习的方法~\cite{microservice-diagnosis-survey}。前者试图显式恢复故障传播结构，并沿依赖关系回溯根因；后者从指标、日志和调用链路中提取判别性表示，直接完成异常定位、排序或分类。两类方法的主要区别在于知识来源不同：前者依赖显式结构，后者依赖数据驱动表示，这种差异决定了它们在不同场景下的适用性。

\subsection{因果推断方法}

这类方法的核心假设是，故障会沿服务调用、部署支撑和资源竞争等依赖关系逐步传播。因此，只要能够较准确地刻画传播结构并识别异常传播方向，就有可能沿传播路径回溯到根因组件~\cite{Pearl2009Causal,Spirtes2000Causation,pham2024causalmicroservicesurvey}。这一假设的有效性依赖于依赖关系的完整性和传播模式的稳定性。

\paragraph{因果图构建与结构发现}
一类工作先构建服务级或实例级依赖图，再结合因果发现或图排序方法恢复异常传播结构。CloudRanger~\cite{Wang2018Cloudranger}和MS-Rank~\cite{Ma2019Ms-rank}从服务依赖关系与监控指标出发构建图结构，并结合因果强度估计或多指标融合对异常传播进行排序。MicroDiag~\cite{Wu2021Microdiag}在组件依赖图基础上结合因果推断，完成更细粒度的性能诊断。ServiceRank~\cite{Ma2021Servicerank}则沿服务依赖图对候选根因进行排序，以提高大规模场景中的定位效率。此类方法共同强调依赖结构在故障传播中的作用，并借助显式图结构增强可解释性。

\paragraph{根因回溯与干预分析}
另一类工作在已有因果图或依赖图上执行根因回溯。MicroRCA~\cite{wu2020microrca}将节点异常程度和边关联强度编码为随机游走权重，并利用 PageRank 式排序~\cite{pagerank-page1999}生成根因候选；REASON~\cite{Wang2023Interdependent}面向相互依赖微服务场景构建跨层级因果网络，并结合带重启的随机游走执行根因定位。CIRCA~\cite{Li2022Causal}将根因定位建模为干预识别问题，通过因果贝叶斯网络中的条件分布变化识别潜在根因指标。近期综述还指出，不同因果发现和图推断方法在真实场景中的效果差异明显，而且对参数设定和数据质量较为敏感~\cite{pham2024causalmicroservicesurvey}。

这类方法的优势是结构清晰、结果可解释，但性能高度依赖图结构质量和因果关系恢复精度。若依赖关系不完整、系统拓扑持续演化，或监控数据存在噪声、稀疏性与概念漂移，沿图传播和干预推理都容易产生系统性偏差~\cite{pham2024causalmicroservicesurvey}。因此，因果推断方法在复杂动态场景下仍表现出较强的结构依赖，这限制了其在快速变化环境中的适用性。

\subsection{特征学习方法}

这类方法将故障诊断视为表征学习和模式识别问题，其核心是从指标、日志和调用链路中学习具有判别力的表示，并利用该表示完成异常检测、根因排序或故障分类。与因果方法相比，这类方法更强调从数据中自动提取特征，对复杂非线性关系具有更强的拟合能力，因而在数据充足时往往表现更好。

\paragraph{单模态特征学习方法}

早期研究多围绕单一模态构建故障表示。日志方向关注从运行日志中捕获异常模式。LogGPT~\cite{loggpt}利用 GPT 式生成模型建模日志序列，并通过强化学习微调策略提升异常检测性能；Landauer 等人~\cite{deep-learning-log-survey}系统梳理了循环神经网络、卷积神经网络和 Transformer 等架构在日志异常检测中的应用。指标方向主要从时序监控数据中检测异常并定位根因，代表工作包括 BARO~\cite{baro-rca}、RobustTAD~\cite{robusttad} 和 DCdetector~\cite{dcdetector}。链路方向则利用调用拓扑与时序关系进行故障定位。以 GAL-MAD~\cite{gal-mad}为例，该方法将微服务调用建模为时空图结构，并结合图注意力网络与长短期记忆网络学习依赖关系。

单模态方法各有优势，也各有明显边界。日志数据语义丰富，但噪声较多，解析质量依赖日志规范；指标数据便于发现数值异常，但难以直接解释故障原因；链路数据能描述传播路径，却难以捕获资源竞争和语义线索。随着这些限制逐渐显现，研究重心开始转向多模态融合。

\paragraph{多模态特征学习方法}

多模态方法同时利用指标中的数值异常、日志中的语义线索和链路中的依赖关系，以形成更完整的故障表示。现有工作大致沿三条路线展开：一是图神经网络的方法，将微服务系统建模为服务依赖图并学习节点表示；二是对比学习方法，通过共享表示提取跨模态不变特征；三是将因果建模与深度学习结合，在表示学习中显式考虑传播关系。

DiagFusion~\cite{zhang2023robust}验证了指标、日志和链路三种模态融合在微服务故障诊断中的有效性；AnoFusion~\cite{anofusion}利用图变换网络建模异构多模态数据间的关联；CHASE~\cite{chase-zhao2024}在多模态图表示学习框架中引入因果超图，以刻画异常在调用路径上的传播关系。TVDiag~\cite{tvdiag}通过统一不同模态的故障视图提高跨场景泛化能力；DynaCausal~\cite{dynacausal}则进一步考虑时序演化和动态依赖。相较于单模态方法，这类方法通常能更充分地利用不同观测信号之间的互补关系。

但多模态特征学习方法仍有局限。其表示往往与特定任务、数据集和故障分布紧密绑定；监督式方法依赖标注样本，无监督方法虽然降低了标注负担，却更偏向异常分离而非经验复用；图神经网络和深度时序模型的内部决策过程不够透明，跨系统迁移时容易受到拓扑差异和模态分布变化的影响~\cite{zhang2023robust,sun2024art}。这种对特定数据分布的依赖限制了方法的泛化能力。

总体来看，传统方法虽然在固定基准和已知故障场景下取得了较好效果，但知识形态仍偏静态。因果推断方法依赖预先构建或离线学习的图结构，特征学习方法依赖围绕特定任务训练出的表示模型。当系统拓扑持续变化或新型故障不断出现时，这些方法往往需要重新建图、重新训练，或反复调整参数。更重要的是，现有表示大多服务于单次定位，难以直接支持相似故障检索与诊断知识绑定，缺乏从历史经验中持续积累的机制。传统方法因此仍缺少一种可独立于具体下游模型的统一故障表征。第三章提出的故障指纹方法正是针对这一缺口展开。
